%%%%%%%% ICML 2026 — Omni-Modal LLM Survey %%%%%%%%%%%%%%%%%

\documentclass{article}

% Recommended packages
\usepackage{microtype}
\usepackage{graphicx}
\usepackage{subcaption}
\usepackage{booktabs}
\usepackage{hyperref}

% Attempt to make hyperref and algorithmic work together better:
\newcommand{\theHalgorithm}{\arabic{algorithm}}

% Blind submission version
\usepackage{icml2026}

\usepackage{amsmath}
\usepackage{amssymb}
\usepackage{mathtools}
\usepackage{amsthm}
\usepackage[capitalize,noabbrev]{cleveref}

% Todonotes for development
\usepackage[textsize=tiny]{todonotes}

% Short running title
\icmltitlerunning{Omni-Modal LLMs: A Comprehensive Survey}

\begin{document}

\twocolumn[
  \icmltitle{Omni-Modal Large Language Models: \\
    A Comprehensive Survey on Architecture, Training, and Evaluation}

  \begin{icmlauthorlist}
    \icmlauthor{Anonymous}{yyy}
  \end{icmlauthorlist}

  \icmlaffiliation{yyy}{Anonymous Institution}

  \icmlkeywords{Omni-modal models, multimodal LLMs, survey, architecture, training, evaluation}

  \vskip 0.3in
]

\printAffiliationsAndNotice{}

% ============================================================
\section{Abstract}
% ============================================================
\begin{abstract}
We present the first comprehensive survey of omni-modal large language models organized around a novel full-lifecycle taxonomy.
While recent years have witnessed an explosion of models capable of jointly processing text, image, audio, and video---exemplified by GPT-4o, Gemini, Qwen3.5-Omni, and ByteDance Seed's Omni---existing surveys focus on narrow aspects such as individual architectural components or training stage decomposition, leaving a fragmented picture of the field.
We bridge this gap by introducing a unified framework that structures the literature along six interconnected dimensions spanning three lifecycle stages: \emph{Build} (Architecture, Training), \emph{Assess} (Capabilities, Evaluation, Safety), and \emph{Apply} (Agents and Deployment).
Through systematic analysis of over 60 recent papers spanning architecture, training, capabilities, evaluation, safety, and agentic applications, we provide structured comparison tables, a unified taxonomy, and an integrated perspective on how architectural choices interact with training strategies, evaluation protocols, and safety mechanisms.
We identify ten critical open challenges---including the lack of truly unified architectures, modality imbalance in training data, the immaturity of omni-modal evaluation benchmarks, and unresolved safety risks across modalities---and propose concrete research directions for each.
We hope this survey serves as both an accessible entry point for newcomers and a structured roadmap for advancing the field.
\end{abstract}

% ============================================================
\section{Introduction}
% ============================================================
\label{sec:intro}

The rapid evolution from uni-modal deep learning models to large language models (LLMs) and subsequently to multimodal large language models (MLLMs) has fundamentally reshaped artificial intelligence research.
Within the span of roughly two years, the frontier has shifted from vision-language models~\cite{liu2024llava,li2023blip2}---primarily handling text and image---to \emph{omni-modal} models that jointly perceive, reason over, and generate across text, image, audio, video, and beyond~\cite{openai2024gpt4o,google2024gemini,qwen2026omni,seed2026omni}.
This transition carries profound implications: a model that can fluidly operate across all sensory modalities mirrors human-like perception more closely than any prior AI system, and unlocks applications ranging from real-time multimodal assistants to embodied agents that see, hear, and act in the physical world.

Yet the research community lacks a unified framework for understanding this rapidly expanding field.
The most closely related survey, ``From Specific-MLLMs to Omni-MLLMs''~\cite{song2024omnimllms}, organizes the literature around four core components (modality encoder, LLM backbone, modality generator, and modality alignment) and a two-stage training pipeline.
While valuable, this component-centric taxonomy overlooks critical cross-cutting dimensions: how training strategies interact with architectural design, how evaluation protocols influence capability assessment, and how safety mechanisms must adapt to the omni-modal setting.
Several other surveys cover adjacent territories---multimodal alignment and fusion~\cite{li2024multimodalalignment}, model architectures~\cite{wu2024evolution}, large multimodal reasoning~\cite{liu2025multimodalreasoning}, and spoken dialogue models~\cite{wang2024wavchat}---but none provide the end-to-end lifecycle perspective that the field now requires.

\textbf{Our contribution.}
This survey bridges the gap by introducing the first \emph{full-lifecycle taxonomy} for omni-modal LLMs.
Our framework organizes the field into six dimensions across three stages:
\begin{enumerate}
    \item \textbf{Build}: How are omni-modal models constructed? We survey \emph{architecture design} (modality tokenization, fusion mechanisms, generation architectures) and \emph{training methodology} (pre-training strategies, modality alignment, efficient training, post-training).
    \item \textbf{Assess}: How do we measure and govern their capabilities? We cover \emph{understanding and generation capabilities}, \emph{evaluation benchmarks}, and \emph{safety and alignment} mechanisms.
    \item \textbf{Apply}: How do these models act in the world? We examine the emerging landscape of \emph{omni-modal agents} and practical deployment considerations.
\end{enumerate}

Specifically, this survey makes the following contributions:
\begin{itemize}
    \item We propose a novel full-lifecycle taxonomy that organizes omni-modal LLM research along Build--Assess--Apply stages, providing a more comprehensive and practically useful framework than existing component-centric surveys.
    \item We systematically review over 60 papers published through mid-2026, capturing the latest developments---including GPT-4o, Gemini 2.0/2.5, Qwen3.5-Omni~\cite{qwen2026omni}, ByteDance Seed's Omni with Context Unrolling~\cite{seed2026omni}, Cosmos~3~\cite{nvidia2026cosmos3}, LatentOmni, DuplexOmni, GuardReasoner-Omni, and emerging omni-agent systems---that postdate prior surveys.
    \item We provide structured comparison tables across all six taxonomy dimensions, enabling researchers to quickly situate new work within the broader landscape and identify under-explored areas.
    \item We identify ten open challenges---ranging from the absence of a unified omni-modal evaluation suite to unresolved cross-modal safety risks---and propose concrete, actionable research directions for each.
\end{itemize}

\textbf{Paper organization.}
Section~\ref{sec:background} establishes preliminaries and contrasts our survey with existing ones.
Section~\ref{sec:taxonomy} presents our full-lifecycle taxonomy in detail.
Sections~\ref{sec:architecture}--\ref{sec:agents} systematically survey each dimension of the framework.
Section~\ref{sec:datasets} catalogs datasets and open resources.
Section~\ref{sec:challenges} identifies open challenges and future directions.
Section~\ref{sec:conclusion} concludes.

% ============================================================
\section{Background and Preliminaries}
% ============================================================
\label{sec:background}

\subsection{From Unimodal to Omni-Modal}

% (Detailed writing to follow in next iteration — tracing the evolution:
%  single modality → vision-language → multi-modal → omni-modal)

The progression of multimodal AI can be characterized by the number and breadth of modalities a single model can handle.
Early multimodal systems operated on fixed modality pairs, typically vision and language~\cite{radford2021clip,li2023blip2,liu2024llava}.
The subsequent wave expanded to three or more modalities by incorporating audio~\cite{ghosal2023speechgpt,deshmukh2023pengi}, video understanding~\cite{lin2023videollava,zhang2023videochat}, or both~\cite{zhan2024anygpt,wu2024nextgpt,su2023pandagpt}.
The current frontier---\emph{omni-modal models}---targets joint perception, reasoning, and generation across text, image, audio, and video within a single unified architecture~\cite{openai2024gpt4o,google2024gemini,qwen2026omni,jiang2024chameleon,seed2026omni}.
Most recently, models such as Qwen3.5-Omni~\cite{qwen2026omni} scale to hundreds of billions of parameters with 256K context and achieve SOTA on 215 audio-visual benchmarks, while ByteDance Seed's Omni~\cite{seed2026omni} introduces \emph{Context Unrolling}---explicit cross-modal reasoning across text, images, video, 3D geometry, and hidden representations before producing predictions.
Concurrently, the community has produced a rapid succession of innovations in efficient tokenization~\cite{placeholder2026latentomni,placeholder2026omnidrop}, quantization~\cite{placeholder2025morphoquant}, real-time full-duplex interaction~\cite{placeholder2026duplexomni}, streaming video understanding~\cite{placeholder2026streamov}, and omni-modal agents~\cite{placeholder2026oscar,placeholder2026sandboxagent,nvidia2026cosmos3}. This explosive growth motivates the need for a structured survey that organizes the literature in a coherent, actionable framework.

\subsection{Key Definitions}

We define the following terms for clarity throughout this survey:
\begin{itemize}
    \item \textbf{Multi-modal model}: A model that processes two or more modalities, typically with separate encoders and a shared representation space.
    \item \textbf{Omni-modal model}: A model capable of jointly processing and/or generating text, image, audio, and video within a unified framework, supporting arbitrary modality combinations as input and output (any-to-any).
    \item \textbf{Modality alignment}: The process of mapping representations from different modalities into a shared embedding space, enabling cross-modal interaction.
    \item \textbf{Any-to-any generation}: The capability to generate content in any supported modality given input in any modality.
\end{itemize}

\subsection{Comparison with Existing Surveys}

\begin{table}[t]
\centering
\caption{Comparison of this survey with existing multimodal surveys. \checkmark{} = covered; $\circ$ = partially covered; $\times$ = not covered.}
\label{tab:survey-comparison}
\begin{tabular}{lcccccc}
\toprule
Survey & Arch. & Train. & Capab. & Eval. & Safety & Agent \\
\midrule
Omni-MLLMs~\cite{song2024omnimllms} & \checkmark & $\circ$ & \checkmark & $\circ$ & $\times$ & $\times$ \\
Architecture Evolution~\cite{wu2024evolution} & \checkmark & $\times$ & $\times$ & $\times$ & $\times$ & $\times$ \\
Alignment \& Fusion~\cite{li2024multimodalalignment} & $\circ$ & \checkmark & $\times$ & $\times$ & $\times$ & $\times$ \\
MLLM Meta-Review~\cite{yang2024surveying} & \checkmark & $\circ$ & \checkmark & \checkmark & $\times$ & $\times$ \\
WavChat~\cite{wang2024wavchat} & $\circ$ & $\times$ & $\circ$ & $\times$ & $\times$ & $\times$ \\
\textbf{Ours} & \checkmark & \checkmark & \checkmark & \checkmark & \checkmark & \checkmark \\
\bottomrule
\end{tabular}
\end{table}

Table~\ref{tab:survey-comparison} systematically compares our survey with existing work.
Our survey is the first to provide integrated coverage across all six dimensions of the omni-modal lifecycle.

% ============================================================
\section{Taxonomy and Framework Overview}
% ============================================================
\label{sec:taxonomy}

% (Framework diagram: Figure 1 — Build → Assess → Apply with 6 dimensions)

We propose a \emph{full-lifecycle taxonomy} that organizes omni-modal LLM research into three sequential stages, each encompassing two core dimensions:

\paragraph{Stage 1: Build.}
This stage addresses the foundational question of \emph{how} to construct an omni-modal model.
\begin{itemize}
    \item \textbf{Architecture} (Section~\ref{sec:architecture}): Modality tokenization, encoder/decoder design, fusion mechanisms, generation architectures, and system-level considerations.
    \item \textbf{Training Methodology} (Section~\ref{sec:training}): Pre-training strategies, modality alignment, efficient training techniques, and post-training refinement.
\end{itemize}

\paragraph{Stage 2: Assess.}
This stage addresses \emph{what} the model can do and \emph{whether} it is safe to deploy.
\begin{itemize}
    \item \textbf{Capabilities} (Section~\ref{sec:capabilities}): Cross-modal understanding, any-to-any generation, and real-time interaction.
    \item \textbf{Evaluation} (Section~\ref{sec:evaluation}): Understanding and generation benchmarks, omni-specific evaluation challenges, and comprehensive assessment frameworks.
    \item \textbf{Safety \& Alignment} (Section~\ref{sec:safety}): Multimodal jailbreak attacks, guardrail mechanisms, and cross-modal value alignment.
\end{itemize}

\paragraph{Stage 3: Apply.}
This stage addresses \emph{where} omni-modal models are deployed and how they interact with the world.
\begin{itemize}
    \item \textbf{Omni-Modal Agents} (Section~\ref{sec:agents}): Agentic omni perception, tool-augmented reasoning, world models, and embodied AI.
\end{itemize}

A sixth cross-cutting dimension---\textbf{Datasets \& Resources} (Section~\ref{sec:datasets})---supports all three stages.

This taxonomy differs fundamentally from prior work: rather than decomposing models by their internal components, we organize the literature by the \emph{lifecycle questions} that researchers and practitioners actually face---how to build, how to assess, and how to deploy omni-modal systems.

% ============================================================
\section{Architecture}
% ============================================================
\label{sec:architecture}

The architecture of an omni-modal LLM must solve a fundamental challenge: how to ingest, represent, fuse, and generate across fundamentally different signal types---discrete text tokens, continuous image pixels, time-varying audio waveforms, and spatiotemporal video frames---within a single unified model.
We organize our survey of architectural paradigms along five dimensions: modality tokenization, encoder/decoder design, fusion mechanisms, generation architectures, and system-level considerations.

\subsection{Modality Tokenization}

Tokenization transforms raw modality signals into a representation suitable for transformer-based processing.
The central trade-off is between modality-specific tokenizers optimized for each signal type and unified tokenizers that map all modalities into a shared vocabulary.

\paragraph{Image tokenization.}
Images are predominantly tokenized via two paradigms.
\emph{Continuous patch embeddings}, popularized by ViT~\cite{dosovitskiy2021vit}, divide an image into fixed-size patches and project each through a learned linear layer, producing a sequence of continuous vectors.
This approach is used by most vision-language models~\cite{liu2024llava,li2023blip2} and remains the default for the vision encoder in omni-modal systems such as GPT-4o and Gemini~\cite{google2024gemini}.
\emph{Discrete visual tokens}, based on vector quantization (e.g., VQ-VAE~\cite{van2017vqvae}, VQGAN), map image patches to indices in a learned codebook, converting images into a ``visual language'' that can be processed by the same autoregressive transformer as text.
Chameleon~\cite{jiang2024chameleon} adopts this approach, tokenizing images into discrete tokens and interleaving them with text tokens for native multimodal generation.

\paragraph{Audio tokenization.}
Audio processing faces a distinct challenge: the high temporal resolution of raw waveforms (typically 16--48 kHz) necessitates aggressive compression.
Two dominant strategies have emerged.
\emph{Spectrogram-based} approaches~\cite{radford2023whisper} convert audio to log-mel spectrograms and process them through a convolutional encoder, producing continuous representations.
\emph{Neural audio codecs} such as EnCodec~\cite{defossez2022encodec} and SoundStream compress audio into discrete tokens at low bitrates (1.5--24 kbps), enabling autoregressive audio generation.
Qwen3.5-Omni~\cite{qwen2026omni} employs a hybrid approach, using separate tokenization pathways for semantic understanding and acoustic generation, bridged by their ARIA alignment mechanism for streaming speech synthesis.

\paragraph{Video tokenization.}
Video introduces the additional dimension of time.
Common strategies include \emph{3D patch extraction} (dividing video into spatiotemporal cubes), \emph{frame-level processing} (treating video as a sequence of images, as in Video-LLaVA~\cite{lin2023videollava}), and \emph{learned temporal aggregation} that compresses frame sequences into compact representations.
Gemini~\cite{google2024gemini} processes video natively through its multimodal encoder, while Qwen3.5-Omni~\cite{qwen2026omni} supports up to 400 seconds of 720P video at 1 FPS via efficient temporal sampling.

\paragraph{Unified tokenization.}
An emerging research direction asks whether a single tokenizer can serve all modalities.
Chameleon~\cite{jiang2024chameleon} demonstrates early success by tokenizing both images and text into a shared discrete vocabulary, enabling mixed-modal autoregressive generation.
ByteDance Seed's Omni~\cite{seed2026omni} pushes further by natively training on text, images, videos, 3D geometry, and \emph{hidden representations}---internal neural states from other models---treating diverse signal types within a unified representational framework.
LatentOmni~\cite{placeholder2026latentomni} explores unified audio-visual latent reasoning, while Conan-embedding-v3~\cite{placeholder2026conan} fuses modality-specific models into omni-modal embeddings.
Unified tokenization promises architectural simplicity but introduces tension between modality-specific compression efficiency and cross-modal vocabulary coherence.

\subsection{Modality Encoders and Decoders}

Beyond tokenization, dedicated encoders extract semantically meaningful features from each modality before passing them to the LLM backbone.
Decoder modules perform the reverse operation for generation.

\paragraph{Encoder architectures.}
Most omni-modal models adopt modality-specific encoders: a ViT-based vision encoder for images~\cite{dosovitskiy2021vit}, a convolutional or Transformer-based audio encoder for speech and sound~\cite{radford2023whisper}, and a video encoder built upon spatial and temporal attention.
The design choices revolve around (1) whether encoders are \emph{frozen} (preserving pre-trained representations, as in BLIP-2's Q-Former~\cite{li2023blip2}) or \emph{fine-tuned} (allowing modality-specific adaptation), and (2) whether encoders are \emph{independent} or \emph{co-trained} with the LLM backbone.
Recent models increasingly favor co-training: Qwen3.5-Omni~\cite{qwen2026omni} jointly optimizes its vision and audio encoders with the core LLM, while Gemini~\cite{google2024gemini} is trained end-to-end as a multimodal model from the start.

\paragraph{Decoder architectures.}
For generation, decoders map latent representations back to modality-specific outputs.
Text generation relies on the LLM's native token decoder.
Image generation typically employs a diffusion-based decoder (e.g., Stable Diffusion-style UNet) or an autoregressive discrete token decoder.
Audio generation uses neural vocoders (e.g., HiFi-GAN) or codec-based decoders that convert discrete audio tokens back to waveforms.
Qwen3.5-Omni's Talker module~\cite{qwen2026omni} incorporates a Hybrid Attention MoE architecture for efficient speech generation, with the ARIA mechanism dynamically synchronizing text and speech units to reduce latency and improve prosody in streaming synthesis.
A key architectural decision is the \emph{generation granularity}: models may generate in a single target modality, support \emph{any-to-any} generation (e.g., NExT-GPT~\cite{wu2024nextgpt}), or produce multiple modalities simultaneously.

\subsection{Fusion Mechanisms}

Modality fusion determines how information from different modalities interacts within the model.
We identify four predominant paradigms:

\paragraph{Cross-attention fusion.}
Cross-attention layers insert modality-specific information into the LLM by attending from text representations to encoded modality features.
BLIP-2's Q-Former~\cite{li2023blip2} exemplifies this approach, using a learnable query transformer to extract visual features relevant to the language model.
This paradigm is lightweight---modality encoders can remain frozen---but limits deep cross-modal interaction to designated attention layers.

\paragraph{Q-Former and Perceiver-based fusion.}
Q-Former and Perceiver architectures compress modality features into a fixed-size set of learned queries, decoupling the modality encoder from the LLM and controlling the number of visual tokens.
This is particularly important for video (where raw frame tokens would exceed context windows) and for efficiency-sensitive deployments.
Flamingo and its successors popularized this approach with Perceiver Resamplers.

\paragraph{Native multimodal attention.}
The most integrated paradigm allows all modality tokens to interact through the same self-attention mechanism.
Chameleon~\cite{jiang2024chameleon} achieves this by tokenizing images into discrete tokens and interleaving them with text in a single sequence, enabling token-level cross-modal attention throughout all transformer layers.
Gemini~\cite{google2024gemini} employs a similar early-fusion design, processing interleaved multimodal sequences natively.
This approach maximizes cross-modal interaction but demands that all modalities share a compatible representation format.

\paragraph{Hybrid attention with modality-specific routing.}
Recent large-scale omni models introduce modality-aware attention mechanisms.
Qwen3.5-Omni~\cite{qwen2026omni} employs a Hybrid Attention Mixture-of-Experts (MoE) framework for both its Thinker (understanding) and Talker (generation) modules, enabling efficient long-sequence inference by routing tokens through modality-specialized expert sub-networks.
ByteDance Seed's Omni~\cite{seed2026omni} introduces \emph{Context Unrolling}: rather than fusing modalities into a single representation, the model explicitly reasons \emph{across} multiple modal representations sequentially, aggregating complementary information from text, images, video, 3D, and hidden states before producing a prediction.
This represents a departure from traditional fusion toward a more deliberative, multi-step cross-modal reasoning process.

\subsection{Generation Architectures}

Omni-modal generation spans text, image, audio, and video outputs, each with distinct architectural requirements.

\paragraph{Autoregressive discrete token generation.}
Models that discretize all modalities into a shared token vocabulary can generate any modality through a unified autoregressive process.
Chameleon~\cite{jiang2024chameleon} and the Gemininative generation pipeline exemplify this approach.
The advantage is architectural uniformity---a single transformer generates all modalities---at the cost of discretization artifacts and large token sequences for high-resolution outputs.

\paragraph{Diffusion-based generation.}
For continuous modalities (images, audio, video), diffusion models dominate generation quality.
Most omni models that generate images or video employ latent diffusion models (e.g., Stable Diffusion) conditioned on text embeddings from the LLM.
This decouples the LLM (which provides semantic conditioning) from the generation backbone (which produces high-fidelity outputs), enabling each component to scale independently.

\paragraph{Hybrid generation pipelines.}
Practical omni models often combine multiple generation paradigms.
GPT-4o generates text and speech autoregressively but uses a separate image generation pipeline.
Qwen3.5-Omni~\cite{qwen2026omni} generates text and speech within its unified Thinker-Talker framework, with the ARIA mechanism ensuring natural prosody and low latency in streaming speech.
The trend is toward tighter integration: reducing the gap between semantic understanding and perceptual generation to enable more natural, responsive multimodal interaction.

\paragraph{Any-to-any generation.}
A defining capability of omni models is supporting arbitrary modality combinations in both input and output.
NExT-GPT~\cite{wu2024nextgpt} and AnyGPT~\cite{zhan2024anygpt} are early examples of any-to-any architectures, connecting modality-specific encoders and decoders through a shared LLM core.
FlowInOne~\cite{placeholder2026flowinone} unifies multimodal generation as image-in, image-out flow matching, while Dynin-Omni~\cite{placeholder2026dyninomni} explores unified diffusion language models for omnimodal generation.
For video and audio, Foley-Omni~\cite{placeholder2026foleyomni} demonstrates unified generation from task-level audio synthesis to complete video soundtracks, and inference-time scaling for joint audio-video generation~\cite{placeholder2026infertimescale} further improves output quality.
Fully unified any-to-any generation---where the same architecture handles all modality pairs without modality-specific modules---remains an open research challenge.

\subsection{System-Level Design}

At scale, omni-modal models face system-level challenges distinct from text-only LLMs.

\paragraph{Mixture-of-Experts for multi-modality.}
MoE architectures~\cite{qwen2026omni} are particularly well-suited to omni models: different modalities naturally exhibit different activation patterns, making sparse expert routing an efficient strategy.
Qwen3.5-Omni's Hybrid Attention MoE enables the model to scale to hundreds of billions of parameters while maintaining tractable inference costs for long multimodal sequences (256K context).

\paragraph{Inference efficiency.}
Omni-modal inference is significantly more expensive than text-only inference due to the computational cost of encoding high-dimensional signals (images, audio, video) and generating in continuous modalities.
Techniques including KV-cache optimization, modality-aware batching, speculative decoding for multimodal tokens, and quantization (e.g., MorphoQuant~\cite{placeholder2025morphoquant} for omni-modal PTQ) are active research areas.
Complementary approaches reduce token counts directly: OmniDrop~\cite{placeholder2026omnidrop} performs layer-wise token pruning guided by query relevance, while Stage-adaptive Token Selection~\cite{placeholder2026stageadaptive} dynamically adjusts token budgets across processing stages.
O-MARC~\cite{placeholder2026omarc} applies memory-augmented compression distillation for efficient video understanding, and StreamOV~\cite{placeholder2026streamov} enables streaming omni-video understanding through evidence-guided memory and response triggering.

\paragraph{Streaming and real-time constraints.}
Real-time omni interaction (e.g., spoken dialogue, live video understanding) imposes strict latency requirements.
Qwen3.5-Omni's ARIA mechanism addresses the encoding efficiency gap between text and speech tokenizers by dynamically aligning text-speech units, achieving natural prosody with minimal latency impact.
Models supporting over 10 hours of audio understanding~\cite{qwen2026omni} require efficient streaming attention mechanisms and intelligent context window management.

\paragraph{Training infrastructure.}
Training omni models at scale requires heterogeneous data pipelines spanning text corpora, image-text pairs, audio-visual content (Qwen3.5-Omni uses over 100 million hours), and multimodal interleaved data.
The coordination of modality-specific pre-training, multi-stage alignment, and joint optimization across different data distributions presents significant engineering challenges that we address in Section~\ref{sec:training}.

% ============================================================
\section{Training Methodology}
% ============================================================
\label{sec:training}

Training omni-modal LLMs presents challenges that go well beyond those of text-only or even vision-language models: heterogeneous data sources spanning fundamentally different signal types, modality-specific convergence dynamics, the risk of catastrophic forgetting across modalities, and the sheer scale of multi-modal datasets.
We organize training strategies along four stages: pre-training, modality alignment, efficient training, and post-training refinement.

\subsection{Pre-training Strategies}

Pre-training an omni-modal model requires reconciling multiple modalities within a single optimization process.
The central challenge is \emph{data heterogeneity}: text corpora, image-text pairs, audio recordings, video streams, and interleaved multimodal documents each follow distinct distributions, tokenization schemes, and information densities.

\paragraph{Data mixture and sampling.}
Determining the optimal ratio of modalities during pre-training is a critical and under-explored design decision.
Most leading omni models train on massive heterogeneous datasets: Qwen3.5-Omni~\cite{qwen2026omni} leverages over 100 million hours of audio-visual content alongside text-vision pairs, while Gemini~\cite{google2024gemini} is trained on a multimodal corpus spanning text, images, audio, and video from the outset.
ByteDance Seed's Omni~\cite{seed2026omni} takes a distinctive approach by including \emph{hidden representations}---internal neural states from pre-trained models---as an additional training modality, enabling the model to leverage complementary information encoded in other networks' latent spaces.
Common strategies include \emph{modality-balanced sampling} (ensuring each modality contributes proportionally to the training objective), \emph{temperature-based mixing} (varying modality ratios over the course of training), and \emph{curriculum learning} (introducing modalities gradually, from simpler to more complex combinations).

\paragraph{Training objectives.}
Omni-modal pre-training typically combines multiple loss functions.
The dominant paradigm is \emph{next-token prediction} over interleaved multimodal sequences, where images, audio, and video are tokenized and concatenated with text tokens into a unified autoregressive stream~\cite{jiang2024chameleon,google2024gemini}.
This is often supplemented with \emph{cross-modal contrastive losses} that pull paired modality representations together in embedding space (inherited from CLIP-style training~\cite{radford2021clip}), and \emph{generation losses} (e.g., diffusion loss for image outputs, vocoder loss for audio) for models that produce continuous modality outputs.
Qwen3.5-Omni~\cite{qwen2026omni} employs separate objective formulations for its Thinker (understanding) and Talker (generation) modules, reflecting the asymmetric complexity of consuming versus producing different modalities.

\paragraph{Training stability and modality imbalance.}
A persistent challenge is \emph{modality competition}: during joint training, the dominant modality (usually text) can suppress learning in weaker modalities.
This manifests as uneven convergence rates across modalities, with text performance improving rapidly while audio or video understanding lags.
Techniques to address this include \emph{modality-specific learning rates}, \emph{gradient scaling} per modality, and \emph{modality dropout} (randomly dropping modalities during training to prevent over-reliance on any single one).
Entrain~\cite{placeholder2026entrain} addresses a related systems challenge: the inherent heterogeneity of multimodal data makes it difficult to balance distributed training workloads across both modalities and batches, requiring modality-aware load balancing.

\subsection{Modality Alignment}

Once pre-trained, omni-modal models must align their modality-specific representations to enable seamless cross-modal interaction.
Alignment strategies can be characterized by \emph{when} and \emph{how} alignment occurs.

\paragraph{Two-stage alignment.}
The most common paradigm, inherited from vision-language models, separates training into a \emph{pre-training stage} (where modality encoders are aligned to the LLM's embedding space) and an \emph{instruction tuning stage} (where the full model is fine-tuned on multimodal instruction-following data).
BLIP-2~\cite{li2023blip2} popularized this with its Q-Former, which is first trained to extract visual features relevant to a frozen LLM, then fine-tuned end-to-end.
This decoupling is computationally efficient but creates a structural bottleneck: the LLM never directly observes raw modality signals, limiting the depth of cross-modal understanding.

\paragraph{End-to-end joint training.}
Increasingly, state-of-the-art omni models adopt end-to-end training from the start.
Gemini~\cite{google2024gemini} and Chameleon~\cite{jiang2024chameleon} are trained as natively multimodal models, with all parameters jointly optimized on interleaved multimodal data.
Qwen3.5-Omni~\cite{qwen2026omni} jointly optimizes vision and audio encoders alongside the core LLM, enabling deeper cross-modal integration at the cost of increased training complexity.
ByteDance Seed's Omni~\cite{seed2026omni} pushes joint training further by co-training on modalities that include not only raw signals but also abstract representations, enabling the model to discover cross-modal correspondences that would be invisible in separate training stages.

\paragraph{Modality-by-modality alignment.}
Some approaches align modalities incrementally: first vision-language, then audio, then video, progressively expanding the model's modality coverage.
While simpler to implement and debug, this risks \emph{catastrophic forgetting}: later modality alignment can overwrite capabilities acquired earlier.
PandaGPT~\cite{su2023pandagpt} and NExT-GPT~\cite{wu2024nextgpt} exemplify modality-by-modality expansion strategies, connecting pre-trained modality-specific modules through a shared LLM core.
SMA~\cite{placeholder2026sma} introduces a \emph{submodular modality aligner} that selects the most informative modality pairs for alignment, improving data efficiency in low-resource settings by going beyond instance-level alignment to leverage the underlying structure of multimodal data.

\subsection{Efficient Training}

Training omni-modal models at scale---hundreds of billions of parameters, millions of hours of multimodal data---demands specialized efficiency techniques.

\paragraph{Modality-aware parallelism.}
Distributed training of omni models introduces challenges beyond standard LLM training.
Different modalities have vastly different computational profiles: encoding a high-resolution image requires orders of magnitude more FLOPs than processing a text token, while video introduces temporal dependencies that complicate batch construction.
Entrain~\cite{placeholder2026entrain} tackles this by chunking and reordering multimodal data to balance workloads across GPUs, accounting for sample-level entanglement between modalities.
Qwen3.5-Omni's Hybrid Attention MoE~\cite{qwen2026omni} enables sparse activation during training: modality-specific tokens are routed through specialized expert sub-networks, reducing the effective computational cost per token while maintaining model capacity.

\paragraph{Parameter-efficient fine-tuning.}
For researchers and practitioners without access to industrial-scale compute, parameter-efficient fine-tuning (PEFT) methods---LoRA, QLoRA, adapter layers---are essential for adapting pre-trained omni models to new tasks or modalities.
However, standard PEFT methods designed for text-only LLMs may not transfer trivially to omni settings, where different modalities benefit from different adapter placements and ranks.

\paragraph{Quantization and compression.}
Post-training quantization (PTQ) for omni models faces unique challenges due to the extreme distribution heterogeneity across modalities.
MorphoQuant~\cite{placeholder2025morphoquant} addresses this with modality-aware quantization, observing that conventional PTQ methods struggle at 4-bit precision for omni-modal LLMs due to disparate outlier patterns across modalities.
Complementary to quantization, token reduction techniques---OmniDrop~\cite{placeholder2026omnidrop} (layer-wise pruning) and Stage-adaptive Token Selection~\cite{placeholder2026stageadaptive} (dynamic token budgets)---reduce the number of tokens processed, directly lowering inference cost without model modification.

\subsection{Post-training}

Post-training refinement---supervised fine-tuning (SFT), reinforcement learning from human feedback (RLHF), and related techniques---is critical for aligning omni models with user intent and improving output quality across modalities.

\paragraph{Staged post-training.}
Recent work on staged post-training for omni-modal models~\cite{placeholder2026stagedposttraining} demonstrates that a carefully sequenced post-training pipeline---instruction tuning followed by visually debiased evaluation and targeted refinement---substantially improves genuine multimodal integration.
A key insight is that standard benchmarks may overstate omni-modal capabilities because visual shortcuts can answer many queries without genuine audio-visual-language integration.
Debiased evaluation reveals that post-training must explicitly enforce cross-modal evidence integration rather than allowing the model to rely on the strongest single modality.

\paragraph{RL post-training and modality collapse.}
Applying standard RL post-training algorithms (e.g., GRPO) to omni models exposes a critical vulnerability: \emph{late-stage modality collapse}~\cite{placeholder2026escapelanguage}.
Because RL algorithms like GRPO apply uniform policy gradients across all tokens---ignoring that different tokens depend unequally on non-text modalities---the optimization process can systematically suppress audio and visual reasoning in favor of text-only shortcuts.
Modality-aware policy optimization, which scales gradients according to each token's dependence on non-text sources, has been proposed as a remedy.

\paragraph{Alignment across modalities.}
Cross-modal alignment in the post-training phase extends beyond instruction following to encompass safety, factuality, and stylistic consistency across generated modalities.
OmniTrace~\cite{placeholder2026omnitrace} introduces a unified framework for generation-time attribution in omni-modal LLMs, identifying which input sources (text, image, audio, video) support each generated statement---a critical capability for verifying factual accuracy in multimodal outputs.
Partial Information Decomposition (PID)~\cite{placeholder2026pid} provides a complementary lens: a decision-level framework that separates the unique, redundant, and synergistic contributions of sensory and linguistic inputs, revealing that much of what appears to be multimodal reasoning may actually be driven by a single dominant modality.
These diagnostic tools are essential for ensuring that post-training genuinely improves omni-modal capabilities rather than masking unimodal shortcuts.

% ============================================================
\section{Capabilities}
% ============================================================
\label{sec:capabilities}

The capabilities of omni-modal LLMs span a wide spectrum---from understanding complex multimodal inputs, to generating content in any target modality, to engaging in fluid real-time interaction.
We structure our survey of capabilities along three axes: cross-modal understanding, any-to-any generation, and real-time interaction.

\subsection{Cross-Modal Understanding}

Cross-modal understanding refers to the ability to jointly reason over information distributed across text, images, audio, and video.
This goes beyond processing each modality independently: the model must integrate complementary signals that may be temporally misaligned, semantically redundant, or subtly contradictory.

\paragraph{Audio-visual-language understanding.}
The most mature capability axis is joint understanding of vision, audio, and text.
Qwen3.5-Omni~\cite{qwen2026omni} achieves state-of-the-art results on 215 audio and audio-visual understanding, reasoning, and interaction benchmarks, surpassing Gemini 3.1 Pro on key audio tasks.
AVI-Bench~\cite{placeholder2026avibench} provides a systematic evaluation of human-like audio-visual intelligence in omni-MLLMs, probing capabilities such as cross-modal grounding, temporal synchronization, and multimodal commonsense reasoning.
ByteDance Seed's Omni~\cite{seed2026omni} demonstrates that training with Context Unrolling---explicit reasoning across multiple modal representations before producing predictions---enables the model to aggregate complementary information across heterogeneous modalities, facilitating a more faithful approximation of shared multimodal knowledge.

\paragraph{Video understanding.}
Video presents unique challenges: extended temporal context (minutes to hours), fine-grained temporal reasoning, and sparse relevant information.
Qwen3.5-Omni~\cite{qwen2026omni} supports over 10 hours of audio understanding and 400 seconds of 720P video at 1 FPS, enabling applications such as long-form video summarization and temporal question answering.
StreamOV~\cite{placeholder2026streamov} tackles streaming video understanding through evidence-guided memory and response triggering, deciding when sufficient information has been accumulated to produce a response rather than processing the entire video.
O-MARC~\cite{placeholder2026omarc} achieves efficient video understanding via memory-augmented compression distillation, compressing long video sequences into compact representations for downstream reasoning.
OmniCap-IF~\cite{placeholder2026omnicapif} benchmarks and improves instruction-following for omni-video captioning, addressing the gap between holistic video understanding and strict adherence to complex user instructions.
VideoOdyssey~\cite{placeholder2026videoodyssey} extends the temporal horizon further, providing a benchmark for ultra-long-context omni-modal video understanding.

\paragraph{Modality interaction and reasoning.}
Understanding \emph{how} modalities interact---rather than merely \emph{that} they can be processed jointly---is an emerging frontier.
Partial Information Decomposition (PID)~\cite{placeholder2026pid} provides a decision-level framework that separates the unique, redundant, and synergistic contributions of different modalities.
Strikingly, PID reveals that much of what appears to be multimodal reasoning may actually be driven by a single dominant modality, with cross-modal synergy contributing far less than expected.
This diagnostic insight is crucial for designing training strategies that genuinely foster cross-modal integration rather than enabling unimodal shortcuts~\cite{placeholder2026stagedposttraining}.
Agentic Active Omni-Modal Perception~\cite{placeholder2026agenticperception} tackles multi-hop audio-visual reasoning, where evidence is sparse, temporally dispersed, and distributed across both audio and visual streams, requiring the model to actively seek out and aggregate information across modalities.

\subsection{Any-to-Any Generation}

A defining promise of omni-modal models is the ability to generate content in any supported modality given input in any modality---\emph{any-to-any generation}.
This capability remains substantially less mature than understanding, but rapid progress is being made across several fronts.

\paragraph{Audio generation.}
Audio generation spans speech synthesis, music, sound effects, and environmental sounds.
Qwen3.5-Omni's Talker module~\cite{qwen2026omni}, powered by the ARIA alignment mechanism, generates streaming speech with human-like emotional nuance across 10 languages, dynamically aligning text and speech units to ensure natural prosody with minimal latency.
The model supports multilingual speech generation and has demonstrated a novel emergent capability termed \emph{Audio-Visual Vibe Coding}---directly performing coding tasks based on audio-visual instructions.

\paragraph{Joint audio-video generation.}
Generating synchronized audio and video is particularly challenging because the two modalities must be temporally aligned at fine granularity.
Foley-Omni~\cite{placeholder2026foleyomni} presents a unified model that spans from task-level audio synthesis (speech, sound effects, music) to complete video soundtrack generation, addressing the real-world need for jointly consistent audio tracks in video production.
Omni-Customizer~\cite{placeholder2026omnicustomizer} tackles end-to-end multimodal customization for joint audio-video generation, enabling the simultaneous preservation of visual identities and vocal timbres across multiple interacting subjects.
Inference-time scaling techniques~\cite{placeholder2026infertimescale} have recently been shown to improve joint audio-video generation quality without additional training, by allocating more computation at inference time to refine both temporal synchronization and semantic alignment.

\paragraph{Long-form video generation.}
Video generation at scale introduces planning challenges distinct from single-frame image generation.
VideoWeaver~\cite{placeholder2026videoweaver} explores agentic long video generation, where an agent framework---leveraging tool use and orchestration capabilities---builds and refines its own video generation workflow, evaluating and evolving generation skills over extended video sequences.
MTAVG-Bench 2.0~\cite{placeholder2026mtavbench} diagnoses failure modes of cinematic expressiveness in multi-talker audio-video generation, revealing specific weaknesses in character consistency, emotional expression, and cross-modal synchronization.

\paragraph{Unified generation architectures.}
Several recent efforts aim to unify generation across modalities under a single architecture.
FlowInOne~\cite{placeholder2026flowinone} reformulates multimodal generation as image-in, image-out flow matching, providing a unified framework for diverse generation tasks.
Dynin-Omni~\cite{placeholder2026dyninomni} explores unified diffusion language models capable of generating across modalities within a single model.
While fully unified any-to-any generation remains an open challenge, these efforts represent important architectural steps toward consolidating what are currently separate generation pipelines.

\subsection{Real-Time Interaction}

Beyond offline understanding and generation, the most transformative omni-modal capability is real-time, full-duplex interaction---the model simultaneously listens, sees, thinks, and speaks, much like a human conversation partner.

\paragraph{Full-duplex multimodal interaction.}
DuplexOmni~\cite{placeholder2026duplexomni} represents the state of the art in real-time multimodal full-duplex interaction, combining seamless real-time perception with complex reasoning and tool use.
Unlike traditional turn-based systems that process one modality at a time, DuplexOmni enables continuous listening, seeing, thinking, and speaking, supporting interruptions, barge-in, and dynamic turn-taking---capabilities essential for natural human-AI interaction.
The key technical challenge is maintaining low end-to-end latency while simultaneously processing multiple modalities and executing tool calls, which DuplexOmni addresses through an optimized streaming architecture.

\paragraph{Streaming speech and low-latency interaction.}
Qwen3.5-Omni's ARIA mechanism~\cite{qwen2026omni} specifically addresses the encoding efficiency gap between text and speech tokenizers that causes instability and unnaturalness in streaming speech synthesis.
By dynamically aligning text and speech units, ARIA achieves stable, natural prosody in conversational speech with minimal latency impact.
The model's ability to understand over 10 hours of continuous audio while responding with sub-second latency demonstrates that long-context understanding and real-time interaction are not mutually exclusive.

\paragraph{Evaluation of real-time capabilities.}
Omni-DuplexEval~\cite{placeholder2026omniduplexeval} provides a dedicated benchmark for evaluating real-time duplex omni-modal interaction, measuring not just response quality but also temporal properties such as latency, turn-taking accuracy, and interruption handling.
OmniPro~\cite{placeholder2026omnipro} complements this with a benchmark for proactive streaming video understanding, where the model must anticipate when to respond based on accumulating visual evidence, rather than passively processing the entire stream.
These evaluation frameworks are critical for moving real-time omni interaction from a qualitative demonstration to a quantitatively measurable capability.

% ============================================================
\section{Evaluation}
% ============================================================
\label{sec:evaluation}

\subsection{Understanding Benchmarks}

\subsection{Generation Benchmarks}

\subsection{Toward Comprehensive Omni Evaluation}

% ============================================================
\section{Safety and Alignment}
% ============================================================
\label{sec:safety}

\subsection{Multimodal Jailbreak and Attacks}

\subsection{Guardrails for Omni-Modal Models}

\subsection{Cross-Modal Value Alignment}

% ============================================================
\section{Omni-Modal Agents}
% ============================================================
\label{sec:agents}

\subsection{Agentic Omni Perception}

\subsection{World Models}

\subsection{Embodied AI}

% ============================================================
\section{Datasets and Resources}
% ============================================================
\label{sec:datasets}

% (Catalog of training datasets, benchmarks, open-source models)

% ============================================================
\section{Challenges and Future Directions}
% ============================================================
\label{sec:challenges}

% (10 open challenges with concrete research directions)

% ============================================================
\section{Conclusion}
% ============================================================
\label{sec:conclusion}

% (Concise summary of contributions and outlook)

% ============================================================
% Acknowledgments
% ============================================================
% (To be added)

% ============================================================
% Bibliography
% ============================================================
\bibliography{main}
\bibliographystyle{icml2026}

\end{document}
